\section{Introduction}
\label{sec:introduction}
\IEEEPARstart{O}{ptical} coherence tomography (OCT) is the standard non invasive way to image the
retina~\cite{huang1991optical}. Every OCT image carries speckle, a multiplicative disturbance that
arises from coherent light scattering inside tissue~\cite{goodman2007speckle}. Speckle blurs the
border between retinal layers, lowers the contrast between tissue and background, and makes
automatic layer segmentation less reliable~\cite{devalla2018drunet}.

Deep denoisers remove much of that disturbance, but they are almost always trained to minimise the
squared error against a clean reference, so a denoiser can raise the peak signal to noise ratio (PSNR)
of a scan while making it less useful to a clinician. Layer boundaries lose definition, the contrast
between tissue and the dark vitreous falls, and thin structures with weak signal fade into the
background. None of these losses is visible in a PSNR number, and none is corrected by training the
same objective harder.

\subsection{Motivation}
The gap we address is not a missing denoiser. It is that nothing sits between the denoiser and the
clinical reader to state in explicit terms what a good OCT image should look like and to repair the
specific ways in which a denoised image falls short. Building that part as another neural block would
hide the reason for every change inside learned weights again, while a clinician, a regulator or a
reviewer who sees a modified scan needs to know which property was judged deficient, where that
judgement was made, and how large a change was allowed. This paper therefore treats post correction as
its own problem. We keep the denoiser frozen and attach a small policy, written as named image
properties, explicit logical rules over them, and explicit limits on the size of any change, that
decides where the image should change and by how much. Every learned map in the decision path is
bounded, and its effect reaches the image only through a named rule.

\subsection{New ideas in this paper}
\label{sec:contributions}
The new ideas are three, and we do not claim that edge filters, gating maps or fuzzy operators are
new on their own.

\begin{enumerate}
\item \textbf{Six clinical properties given in closed form.}
Section~\ref{sec:predicates} defines six predicates that need no ground truth. Each returns a score
for the whole image and a map of where the property is unsatisfied. Both are formulas, so the input
to the rule layer can be recomputed by anyone from the released code.

\item \textbf{A fuzzy allocation rule with a proved and measured sensitivity bound.}
Section~\ref{sec:negotiator} builds the allocation from Lukasiewicz connectives.
Section~\ref{sec:stability} defines what stability means here and proves a Lipschitz bound, and
Section~\ref{sec:theory_check} measures the true sensitivity on real images and compares it with the
bound. A bound that is never checked against data is a claim, not a guarantee.

\item \textbf{A safety layer that decides how much of a proposal to keep.}
Section~\ref{sec:safety} states three constraints at three scales and shows how the number of
satisfied constraints sets the fraction of the proposed change that is applied. A margin condition
under which the correction cannot lower the contrast to noise ratio is proved and then tested
numerically.
\end{enumerate}

